\subsection{What the Lean development proves}

The accompanying Lean project pins Lean~4.19.0 and Mathlib revision
\begin{center}\small\nolinkurl{c44e0c8ee63ca166450922a373c7409c5d26b00b}.\end{center}
It has $32$ proof modules. A separate audit module enumerates
all $2068$ declarations whose defining module is one of them, including
private and generated declarations, and computes their transitive axioms.
Of these, $1319$ are auxiliary declarations that the Lean~4.19 code
generator stores for executable definitions; beyond the standard axioms
they may refer only to its placeholder \texttt{lcProof} and to one another,
and no other declaration depends on them through an axiom. Each of the other $749$
declarations is a theorem, definition or inductive construction, not an
axiom, and depends only on
\[
 \{\texttt{propext},\ \texttt{Classical.choice},\ \texttt{Quot.sound}\}.
\]
There is no \texttt{sorry}, no added axiom and no native evaluation;
in particular \texttt{Lean.ofReduceBool} does not occur.

The development defines the inversion-weighted polynomial $P_q$ of an
ordered matrix by its permutation sum, exactly as in
Section~\ref{sec:endpoint}, and proves that $D_q$ is its derivative.
It proves, for every real matrix $V$ with three columns, the identity
$2D_1(VV^T)=\binom n2\|f\|_{\F}^2-\|G\|_{\F}^2$ of
Lemma~\ref{lem:endpoint}, with both Fischer norms written as
factorial-weighted coefficient sums of the actual product and wedge
polynomials. It proves the four-row recurrence of
Lemma~\ref{lem:recurrence}, the parity decomposition
\eqref{eq:certP}--\eqref{eq:certT}, the semantics of an executable
sparse integer implementation of the recurrence, and certificate
soundness: for every list $L$ of integer types, if the integer test
$T_U>\binom n2P$ computed from $L$ succeeds, then the Gram matrix of the
row word of $L$ has $D_1<0$, and there are a rational matrix of the same
order that is positive definite and rationals $0<q_1<q_2<1$ with
$P_{q_2}<P_{q_1}$. The last step is a general reduction: a real
positive-semidefinite matrix with negative endpoint derivative yields a
rational positive-definite matrix and a rational interior decreasing
pair. Identity padding is proved for the actual polynomial, together
with preservation of positive definiteness and rationality; one fixed
rational pair therefore serves every larger order.

For the frozen input, the development defines the $\WitnessN$ ordered rows
from the reversed type list, proves that the Gram matrix is positive
semidefinite, and proves that the row matrix and the Gram matrix have
rank exactly three. The rank proof uses the rows with zero-based indices
$0,1,894$, whose matrix is
\[
 \begin{pmatrix}-3&2&20\\-3&-2&20\\3&2&20\end{pmatrix},
 \qquad \det=480,
\]
through a lemma about the outermost type that avoids expanding all
$224$ nested row constructors. This is a different minor from the one
in Table~\ref{tab:witness}; both certify the same rank. Finally it
proves that each of the following hypotheses implies the all-orders
statement
\[
 \exists q_1,q_2\in\mathbb Q,\ 0<q_1<q_2<1,\quad
 \forall N\ge\WitnessN\ \exists A_N\in\mathbb Q^{N\times N},\quad
 A_N\succ0,\quad P_{q_2}(A_N)<P_{q_1}(A_N):
\]
the integer test above for the frozen type list, and the Fischer-norm
inequality $\binom{\WitnessN}2\|f\|_{\F}^2<\|G\|_{\F}^2$ for the frozen rows.

The project also contains scalar lemmas used by the asymptotic
sections: bounds for the spherical Gini gain, logarithmic and radial
errors, rational score margins and condition-number arithmetic for the
families of Theorems~\ref{mf:20002}, \ref{mf:257} and~\ref{rf:228}, and
two finite matrix steps of Section~\ref{sp:section}. These lemmas
formalize the indicated inequalities only; they do not formalize the
probabilistic constructions or limits that supply their hypotheses.

\subsection{What is checked outside Lean}

The evaluation of the integer test on the frozen input is not carried
out in Lean. Kernel reduction of the $224$-step recurrence, whose
polynomials have up to about $10^5$ coefficients of several thousand
digits, is far beyond practical limits. An earlier stage of this work
evaluated it with Lean's native code generator through
\texttt{native\_decide}, which adds \texttt{Lean.ofReduceBool} and hence
the compiler and runtime to the trusted base \cite{LeanNativeDocumentation};
that module is not part of the present development. Instead, the
hypothesis of the conditional theorems is checked by the two exact
programs of Section~\ref{sec:certificate}: the integer implementation of
the four-row recurrence, and the independent row-by-row modular
computation, which determines $P$, $T_U$ and $T$ through the Chinese
remainder theorem. The finite counterexample is therefore a formal
reduction together with an exact computation; it is not a closed
Lean theorem.

The explicit matrix $A=MB+I_n$ and the explicit pair $q_1=1-2/Q$,
$q_2=1-1/Q$ of Section~\ref{fr:section} have conventional proofs; the
formal positive-definite conclusion is existential. The mean-field
endpoint functional, Gaussian concentration, subcritical asymptotic,
spherical saddle analysis, algebraic grids, operator factorization,
spectral obstruction, random-order theorem and coherent-ordering
theorem are conventional proofs as well. Their finite matrix and
scalar ingredients are formalized only where stated above.
